% Subsection 5.3.1: LoRA 有效性分析

本节聚焦训练侧主效应，基于 5.2.1 节表~\ref{tab:ch5-direct} 取出 $\{$H4, H7, H9$\} - $H1 三条差分量，从「引入 LoRA 是否有效」「CoT 监督是否进一步有效」「训练侧 few-shot 是否有独立增益」三层依次展开。

\subsubsection{纯代码 LoRA 的基础增益}
令 $\Delta_1 = \mathrm{H4} - \mathrm{H1}$。在 Direct 推理下，$\Delta_1$ 直接衡量把「纯代码 solution」作为监督目标对基座带来的增益。预期结论分三种情形：
\begin{itemize}
    \item $\Delta_1 > 0$ 且在 pass@1 与 pass@10 上同号：说明 LoRA 提升的是「稳定产出正确代码」的能力，而非单纯扩大候选多样性；
    \item $\Delta_1 > 0$ 但仅在 pass@1 上显著：说明训练主要收敛到「更高概率给出主流解答」，对长尾题的覆盖改善有限；
    \item $\Delta_1 \approx 0$ 或 $< 0$：说明训练集与评测集的分布差异主导了结果，或 LoRA 超参未匹配本任务——此时应回到第四章的超参数设置做敏感性检查。
\end{itemize}
实测结果：$\Delta_1 = +8.1$\,pp（pass@1）、$+6.7$\,pp（pass@5）、$+6.1$\,pp（pass@10），三档同号正向且 pass@1 与 pass@10 同时显著（pass@1 增益最大），落入第一种情形——\textbf{LoRA 提升的是「稳定产出正确代码」的能力，而非单纯扩大候选多样性}。增益落在「同分布 SFT 合理区间」内（既显著大于 3\,pp 的「LoRA 欠拟合」下界，也明显低于 15\,pp 的「训练-评测题面重叠 / 数据泄漏」上界），结合 4.5 节训练损失曲线收敛终点（约 0.18）可进一步排除 target\_modules 覆盖不全或学习率失配等结构性故障。

\subsubsection{CoT 监督的增量}
令 $\Delta_2 = \mathrm{H7} - \mathrm{H4}$。该量固定推理侧为 Direct，只改变训练侧的监督目标（纯代码 $\to$ 推理+代码），因此严格剥离了「训练时是否见过 CoT 格式」的独立贡献。预期以下两条表现之一：
\begin{itemize}
    \item $\Delta_2 > 0$：即使推理时仍走单阶段 Direct，模型也能从训练阶段吸收推理习惯，说明 CoT 监督把「思考链」内化进参数；
    \item $\Delta_2 \le 0$：监督目标中的推理段对 Direct 推理反而造成干扰，说明 CoT 知识必须在推理侧显式展开才有收益，这一结论将与 5.3.2 节「训练-推理对称性」讨论互为印证。
\end{itemize}
实测结果：$\Delta_2 = -3.3$\,pp（pass@1）、$-3.0$\,pp（pass@5）、$-2.5$\,pp（pass@10），三档同号负向，落入第二种情形——\textbf{CoT 知识必须在推理侧显式展开才有收益，「把推理监督进参数」对 Direct 单阶段推理反而构成格式开销}。结合 5.3.2 节交互项 $(H8-H7)-(H2-H1)=+5.0$\,pp 的强正向 DiD，可进一步推断：CoT 监督的价值主要体现在「同分布推理」场景下的协同放大，而非单方向地内化为 Direct 推理可调用的参数知识。

\subsubsection{训练侧 few-shot 的独立增益}
令 $\Delta_3 = \mathrm{H9} - \mathrm{H7}$。两组的训练目标都是「推理+代码」，差别仅在训练 prompt 是否前置 $k$ 条示例。该差分回答本文区别于 8 组版本析因的核心问题——\textbf{在监督目标已经含 CoT 的前提下，继续把 few-shot 写进训练上下文是否还有效}。三种典型情形：
\begin{itemize}
    \item $\Delta_3 > 0$：训练时见过 few-shot 会让 LoRA 学到「如何使用示例」本身，该信号会在推理侧 Direct 下仍然起作用；
    \item $\Delta_3 \approx 0$：few-shot 的作用局限在推理侧，训练阶段注入示例仅起到「扩大上下文规模」的效果，没有形成可内化的参数知识；
    \item $\Delta_3 < 0$：训练 prompt 过长导致 assistant 段在 \texttt{MAX\_LENGTH=1024} 约束下被截断，反而削弱监督信号——此时应增大 \texttt{MAX\_LENGTH} 或降低 $k$ 再做敏感性对照。
\end{itemize}
实测结果：$\Delta_3 = +1.3$\,pp（pass@1）、$+1.1$\,pp（pass@5）、$+0.9$\,pp（pass@10），三档同号弱正向，落入第一种情形——\textbf{训练时见过 few-shot 会让 LoRA 学到「如何使用示例」本身，该信号在推理侧 Direct 下仍能起作用，但量级远小于 $\Delta_1$}。该结果与 5.3.2 节 few-shot 边际一致性判定（基座 $H3-H2=+4.2$\,pp 与 \texttt{LoRA\_code} 上 $H6-H5=+3.7$\,pp 量级一致）互相印证：few-shot 既存在推理侧的结构性边际，也存在训练侧的弱内化效应，两条路径量级差约一个数量级。

\subsubsection{LoRA 与基线的稳定性对比}
除 pass@$k$ 主指标外，本节还关注 LoRA 对「生成稳定性」的改变。稳定性由两项观测量间接刻画：原始 \texttt{solution} 自检率的偏离幅度（环境一致性）与 CoT 两阶段生成的超时率（运行时稳定性）。实测结果：
\begin{itemize}
    \item \textbf{原始 solution 自检率}：H1/H4/H7/H9 四组 Direct 评测 JSON 中 \texttt{original\_solution\_pass\_rate} 均为 99.60\%，三组 LoRA 与基座完全一致 $\Rightarrow$ 执行器与依赖环境稳定，pass@$k$ 差分可纯归因于方法侧；
    \item \textbf{CoT 端超时率}：在 6 组 CoT 评测 \texttt{generation\_stats.timeouts} 中，CoT 监督 LoRA 系列（H8 = 31、H10 = 34，对应 3.1\%、3.4\%）反而\textbf{低于}基座 CoT（H2 = 38、H3 = 41，对应 3.8\%、4.1\%），说明 CoT 监督让模型学会更克制地组织推理段长度，而非把训练过拟合到少数长解答模板上——这一现象与「CoT 监督把推理预算的隐式控制内化进参数」的解释一致。
\end{itemize}
两项稳定性指标在 LoRA 三档上均未出现劣化，进一步支撑 $\Delta_1$ 主效应的可靠性。